\section{Introduction}
\label{sec:intro}

Film production is a multi-million dollar industry where shooting schedules directly impact budgets.
The \emph{Movie Scene Scheduling Problem} (MSSP) determines the order and day assignments of scenes to minimize total production cost---comprising actor wages, hierarchical location transfer costs, deadline penalties, and criticality-weighted resource utilization---subject to a rich set of real-world constraints: actor availability calendars, a three-tier geographic location hierarchy (city--town--location), crew travel logistics that restrict town changes across consecutive days, daily shooting capacity, precedence relations, time windows, and labor regulations limiting consecutive working days.

The MSSP is NP-hard, reducing to a variant of the resource-constrained project scheduling problem (RCPSP)~\citep{brucker1999}.
Exact methods based on integer linear programming (ILP) can find provably optimal solutions for small instances (10--15 scenes) but become intractable as problem size grows~\citep{liu2019, tekin2023}.
Metaheuristic approaches---including tabu search~\citep{liu2019}, particle swarm optimization~\citep{long2020, tekin2023}, and ant colony optimization~\citep{long2020}---offer scalable alternatives but operate on flat scene permutations, ignoring the structural coupling between scenes that share actors or locations.
Moreover, existing formulations employ simplified constraint models that omit critical production logistics such as hierarchical location structures, town continuity requirements, and actor labor regulations.

We observe that real film schedules exhibit \emph{shooting blocks}: groups of scenes that are tightly coupled through shared actors and co-located sets, and are typically filmed consecutively in practice.
This observation motivates a decomposition approach that exploits this structure.

\subsection{Contributions}

We propose \textbf{BDM-AOS} (Bi-level Decomposition Matheuristic with Adaptive Operator Selection), which makes the following contributions:

\begin{enumerate}
    \item \textbf{Comprehensive MILP formulation:} We present the most complete MSSP formulation to date, with 13 constraint families over a three-tier geographic hierarchy (city--town--location), actor--scene--day and location--scene--day linking variables, town-uniqueness and consecutive-day town-continuity constraints, criticality-weighted objective terms, and actor consecutive working-day limits.
    This extends prior formulations~\citep{liu2019, tekin2023} with geographic logistics constraints and resource criticality modeling.

    \item \textbf{Scene Interaction Graph (SIG):} A novel graph-theoretic representation where nodes are scenes and edge weights quantify coupling through actor sharing (Jaccard similarity), location sharing, and transfer cost proximity. Louvain community detection~\citep{blondel2008} on the SIG partitions scenes into shooting blocks, dramatically reducing the combinatorial search space.

    \item \textbf{Bi-level Decomposition with Greedy Reassembly:} The upper level uses NSGA-II~\citep{deb2002} to evolve permutations of shooting blocks, simultaneously optimizing total cost and makespan. The lower level applies a global greedy decoder that packs scenes into days in block-guided order, enforcing all 13 constraint families. This is the first decomposition-based approach for the MSSP.

    \item \textbf{Q-Learning Adaptive Operator Selection (AOS):} A reinforcement learning mechanism that dynamically selects among a pool of six crossover-mutation combinations, including two novel problem-specific operators:
    \begin{itemize}
        \item \emph{Location-Aware Crossover (LAX):} preserves sub-sequences where consecutive blocks share the same primary location;
        \item \emph{Actor-Continuity Mutation (ACM):} reduces actor idle gaps by repositioning blocks to improve actor schedule continuity.
    \end{itemize}

    \item \textbf{Multi-objective formulation:} We present the first bi-objective MSSP formulation optimizing cost and makespan simultaneously, providing a Pareto front of trade-off solutions.

    \item \textbf{Experimental evaluation:} Comprehensive experiments on 8 instance sizes (10--100 scenes) with 30 independent seeds per algorithm, five ablation variants, hypervolume and inverted generational distance quality indicators, and Friedman tests with Holm's post-hoc correction~\citep{garcia2009, demsar2006, holm1979}.
\end{enumerate}

The remainder of this paper is organized as follows.
Section~\ref{sec:related} reviews related work.
Section~\ref{sec:methodology} presents the problem formulation and the BDM-AOS algorithm.
Section~\ref{sec:experiments} describes the experimental setup and results.
Section~\ref{sec:discussion} discusses findings and implications.
Section~\ref{sec:limitations} addresses limitations.
Section~\ref{sec:conclusion} concludes with future research directions.
